\documentclass[justified, nobib,a4paper]{tufte-handout}

\title{Different forms of the energy equation}

\author[Wojciech Sadowski]{Wojciech Sadowski}

%\date{28 March 2010} % without \date command, current date is supplied

% \geometry{showframe} % display margins for debugging page layout

\usepackage{fontspec}
% \setmainfont[Renderer=Basic, Numbers=OldStyle, Scale = 1.0]{TeX Gyre Pagella}
% \setsansfont[Renderer=Basic, Scale=0.90]{TeX Gyre Heros}
% \setmonofont[Renderer=Basic]{TeX Gyre Cursor}

% Set up the spacing using fontspec features
\renewcommand\allcapsspacing[1]{{\addfontfeature{LetterSpace=15}#1}}
\renewcommand\smallcapsspacing[1]{{\addfontfeature{LetterSpace=10}#1}}

% setting ptions for included pictures
\usepackage{graphicx} % allow embedded images
\graphicspath{{graphics/}} % set of paths to search for images
\setkeys{Gin}{
  width=\linewidth,
  totalheight=\textheight,
  keepaspectratio
}

\hypersetup{colorlinks}

\usepackage[compatibility=false, font=footnotesize]{caption}
% * options are necessary to fix the bug that autoref was causing,
%   figure was constantly changed to section
\DeclareCaptionFont{blue}{\color{blue}}
\captionsetup{labelfont={blue}}

% color and color palettes
\usepackage{xcolor}

% drawings collor palette
\definecolor{p11}{HTML}{33cfff}
\definecolor{p12}{HTML}{ffa733}
\definecolor{p13}{HTML}{ff4133}


% drawing inside tex
\usepackage{tikz} 
\usetikzlibrary{decorations.markings}
\usetikzlibrary{calc}
\usetikzlibrary{shapes}

\usepackage{pgfplots}
\pgfplotsset{compat=newest}
\usepgfplotslibrary{fillbetween}
\usepackage{xcolor}

% math things
\usepackage{
  amsmath,  % extended mathematics
  amsthm,
  siunitx,  % non-stacked fractions and better unit spacing
  bm,       % bold math symbols
  physics2,  % a set of usefull symbols for differentiation 
  fixdif,
  derivative
}

\usephysicsmodule{ab, nabla.legacy}

% table things
\usepackage{
  booktabs, % book-quality tables
  multicol  % multiple column layout facilities
} 


% writing references
\usepackage{cleveref}  % must be loded after amsmath

% quoting text and friends
\usepackage{fancyvrb}         % extended verbatim environments
\fvset{fontsize=\normalsize}  % default font size for fancy-verbatim environments
\usepackage{
  dirtytalk, % provides \say command for easy quiting 
  nth       % provides \nth command for 1-st, 2-nd, etc
}        


% citing sources
\usepackage{natbib}
\setcitestyle{authoryear}
\bibliographystyle{plainnat}

% misc 
\usepackage{lipsum}   % filler text

\tikzset{
  set arrow inside/.code={\pgfqkeys{/tikz/arrow inside}{#1}},
  set arrow inside={end/.initial=>, opt/.initial=},
  /pgf/decoration/Mark/.style={
    mark/.expanded=at position #1 with
    {
      \noexpand\arrow[\pgfkeysvalueof{/tikz/arrow inside/opt}]{\pgfkeysvalueof{/tikz/arrow inside/end}}
    }
  },
  arrow inside/.style 2 args={
    set arrow inside={#1},
    postaction={
      decorate,decoration={
        markings,Mark/.list={#2}
      }
    }
  },
}

\input{../src/commands.tex}
\input{../src/symbols.tex}
\newtheorem{theorem}{Theorem}
\newtheorem{lemma}{Lemma}

\newcolumntype{L}[2]{>{\raggedleft#2}p{#1\textwidth}}


\newcommand{\Q}[1]{\subsection{Q:~{#1}}}


\begin{document}

\maketitle% this prints the handout title, author, and date

\tableofcontents

\section{Conservation of mass and momentum}
Mass conservation law can be expressed as\footnote{%
  Meaning of the symbols
  \begin{description}
    \item[\(\rho\)] density
    \item[\(\bm{u}\)] velocity vector
    \item[\(u_i\)] \(i\)-th component of \(\bm{u}\)
    \item[\(\bm{x}\)] spatial coordinate
    \item[\(t\)] time
  \end{description}
}
\begin{equation}
  \pdv{\rho}{t} + \pdv{\rho u_i}{x_i} = 0.
  \label{eq:mass-conserv}
\end{equation}
The first term denotes the change of density in time in an infinitesimal volume
element and the second describes the change of density due to the flux of mass
(\(\rho\bm{u}\)) through the boundaries of the volume element. 

Momentum equation can be expressed as\footnote{%
  Meaning of the symbols
  \begin{description}
    \item[\(p\)] pressure
    \item[\(\bm{\sigma}\)] viscous stress tensor
    \item[\(\bm{f}\)] body force (e.g., gravity)
  \end{description}
}
\begin{equation}
  \pdv{\rho u_i}{t} + \pdv{\rho u_i u_j}{x_j} = -\pdv{p}{x_i} + \pdv{\sigma_{ij}}{x_j} + \rho f_i.
  \label{eq:momentum-conserv}
\end{equation}

\subsection{Material derivative}
The material derivative \(\mdv{}/{t}\) of a variable \(f(\vec{x}, t)\) convected with the
fluid, is defined as 
\[
  \mdv{f}{t} = \underbrace{\pdv{f}{t}}_{I} + \underbrace{u_i \pdv{f}{x_i}}_{II}
\]
where the first term (\(I\)) describes the observed change of \(f\) due to
\emph{real} change of \(f\) with time, whereas the second term denotes the
change due to the different value of \(f\) convected with the flow field.
Consider the temperature rise as illustrated in \cref{fig:material}. The
temperature observed in a point \(\vec{x}\) in the lake with unmoving water can
rise due to the heating of the sun. This would correspond to \(\pdv{T}/{t}\).
If the water moves and the water upstream is hotter, then it will also be ob
observed as the rise of temperature \(T\) in point \(\vec{x}\), which is
expressed as \(u_i\pdv{T}/{x_i}\).

\begin{marginfigure}
  \begin{tikzpicture}
    \draw[thick, fill=blue!20] (0, 0) -- (4, 0) -- (4, 2) -- (0, 2) -- cycle;
    \node at (0, 0) [above, anchor=south west] {lake with standing water};
    \filldraw [red] (1, 1) circle (2pt) node [anchor=south west] {observer sees $T\uparrow$};
    \draw [stealth-,orange,thick] (2,1.7) -- (3,2.5) node {sun rays};
  \end{tikzpicture}
  \begin{tikzpicture}
    \draw[thick, fill=blue!20] (0, 0) -- (4, 0) -- (4, 2) -- (0, 2) -- cycle;
    \node at (0, 0) [above, anchor=south west] {river with current $u$ and $T_l > T_r$};
    \node[red] at (0, 1) [right] {$T_l$};
    \node[red] at (4, 1) [left] {$T_r$};
    \draw[thick, -stealth] (2.5, 0.5) -- (3.7, 0.5) node [above] {$u$};
    \filldraw [red] (1, 1) circle (2pt) node [anchor=south west] {observer sees $T\uparrow$};
  \end{tikzpicture}
  \caption{Illustration of the meaning of terms in material derivative. Upper schematic for \(I\) and 
  lower for \(II\).}\label{fig:material}
\end{marginfigure}

We want to use the following formula to convert between the conservative and non-conservative 
formulations of the equations:
\[
  \rho \mdv{f}{t} 
  = \rho \ab(\pdv{f}{t} + u_i \pdv{f}{x_i})
  = \pdv{\rho f}{t} + \pdv{\rho u_i f}{x_i}.
\]
We can quickly prove that the forms are equivalent using the product rule and continuity equation
\[
  \pdv{\rho f}{t} + \pdv{\rho u_i f}{x_i} 
  = \rho\pdv{f}{t} + \pdv{\rho}{t}f + \rho u_i \pdv{f}{x_i} + \pdv{\rho u_i}{x_i}f=
  \rho\mdv{f}{t} + f\ab(\pdv{\rho}{t} + \pdv{\rho u_i}{x_i})
\]
The last term on the right is of course 0 due to \cref{eq:mass-conserv}.

\section{Kinetic energy equation}
Kinetic energy \(K\) is defined as\footnote{%
  The operation of multiplication of two vectors (variables with one index)
  with the same index repeated is a dot product, i.e., \(u_i u_i \equiv \bm{u}\bm{\cdot}\bm{u}\).
}
\[
  K = \frac{1}{2}u_i u_i.
\]
If we take any derivative of \(K\), here written as \(\pdv{K}/{s}\), with \(s = x_i\) or \(t\), 
we can notice something interesting:
\[
  \pdv{K(u_i(s, \ldots))}{s} = \pdv{K}{u_i}\pdv{u_i}{s} = \pdv{}{u_i}\ab(\frac{1}{2}u_i^2)\pdv{u_i}{s} = u_i \pdv{u_i}{s}.
\]
We just computed the derivative using the chain rule,\footnote{
  \(\odv{f(g(x))}{x} = \odv{f}{g}\odv{g}{x}\)
}
and we will use the same method to derive the equation for \(K\). If we
multiply \cref{eq:momentum-conserv} with \(u_i\) 
\[
  u_i\pdv{\rho u_i}{t} + u_i\pdv{\rho u_i u_j}{x_j} = -u_i\pdv{p}{x_i} + u_i\pdv{\sigma_{ij}}{x_j} + \rho u_i f_i,
\]
we will lead to the same situation as above, i.e., a derivative of \(u_i\)
times \(u_i\) itself. This can be simplified now, by working backwards:
\begin{equation}
  \begin{split}
    u_i\pdv{\rho u_i}{t} + u_i\pdv{\rho u_i u_j}{x_j} &= 
    u_i u_i\pdv{\rho}{t} + \rho u_i \pdv{u_i}{t} 
    + \rho u_j u_i \pdv{u_i}{x_j} + u_i u_i \pdv{\rho u_j}{x_j} \\
    &= u_i u_i \underbrace{\ab( \pdv{\rho}{t} + \pdv{\rho u_j}{x_j})}_{=0} + 
    \rho u_i \pdv{u_i}{t} + \rho u_j u_i \pdv{u_i}{x_j}
  \end{split}
  \label{eq:k1}
\end{equation}
In the last two terms we are using the chain rule as before
\[
  \rho u_i \pdv{u_i}{t} + \rho u_j u_i \pdv{u_i}{x_j} 
  = \rho \pdv{}{t}\ab(\frac{1}{2}u_i u_i) 
  + \rho u_j \pdv{}{x_j}\ab(\frac{1}{2}u_i u_i)
  = \rho \pdv{K}{t} + \rho u_j \pdv{K}{x_j}
\]
Finally we use a standard trick of changing between the primitive and conservative 
form, i.e., we rewrite each term as
\[
  \rho \pdv{K}{t} = \pdv{\rho K}{t} - \pdv{\rho}{t} K,
\]
\[
  \rho u_j \pdv{K}{x_j} = \pdv{\rho u_j K}{x_j} - \pdv{\rho u_j} K,
\]
and substitute into \cref{eq:k1} using \cref{eq:mass-conserv}to simplify:
\[
 \rho u_i \pdv{u_i}{t} + \rho u_j u_i \pdv{u_i}{x_j} = 
  \pdv{\rho K}{t} + \pdv{\rho u_j K}{x_j} - \underbrace{\ab(\pdv{\rho}{t} + \pdv{\rho u_j}{x_j})}_{=0} K
  = \pdv{\rho K}{t} + \pdv{\rho u_j K}{x_j}
\]
Leading to the conservation equation for kinetic energy
\[
  \pdv{\rho K}{t} + \pdv{\rho u_j K}{x_j} 
  = -u_i\pdv{p}{x_i} + u_i\pdv{\sigma_{ij}}{x_j} + \rho u_i f_i.
\]

\section{Total energy equation}
\[
  \odv{}{t}\ab(\text{Energy in the control volume } \Omega) = \text{Power}
\]
We will call power \(N\) or rate of work done per unit time. There can be
different sources of \(N\):
\begin{enumerate}
  \item work of body forces\footnote{%
      Meaning of symbols:
        \begin{description}
          \item[\(\Omega\)] arbitrary control volume
          \item[\(S\)] surface of the volume
          \item[\(\vec{n}\), \(n_i\)] normal of the surface
          \item[\(\vec{q}\), \(q_i\)] heat flux
          \item[T] temperature
        \end{description}
      } 
    \[
      N = \int\limits_\Omega\rho f_i u_i \d \Omega
    \]
  \item work of surface forces on the boundary of the control volume
    \[
      N = \oint\limits_S u_i\ab(-p\delta_{ij} + \sigma_{ij})n_j \d S \stackrel{\text{Green th.}}{=}
      \int\limits_\Omega \pdv{}{x_j}\ab[u_i\ab(-p\delta_{ij} + \sigma_{ij})] \d \Omega
    \]
  \item heat flux through the surface (can be according to Fick's law, i.e.,
    \(\vec{q} = -\lambda\grad{T}\))
    \[
      N = \oint\limits_S - q_j n_j \d S = \int\limits_\Omega -\pdv{q_j}{x_j}\d \Omega 
    \]
  \item balance of local point sources (e.g., due to chemical reactions)
    \[
      N = \int\limits_\Omega \rho Q \d \Omega
    \]
\end{enumerate}
The total energy in \(Omega\) is given as 
\[
  \int\limits_\Omega \rho E \d \Omega 
  = \int\limits_\Omega \rho e + \frac{\rho u_i u_i}{2} \d \Omega 
  = \int\limits_\Omega \rho e + \rho K \d \Omega 
\]
We need to simplify this expression 
\[
  \odv{}{t}\ab(\text{Energy in the control volume } \Omega) = 
  \odv{}{t}\int\limits_\Omega \rho E \d \Omega,
\]
and the problem is that \(\Omega\) can be potentially deformed by the fluid flow. 
We have to use Reynolds Transport theorem given as stated below to achieve
this.

\begin{theorem}[Reynolds transport theorem, RTT]
	Given:
	\begin{enumerate}[(i)]
		\item a material volume \(\Omega\), whose shape and position are varying with time,
		      i.e.\ \(\Omega = \Omega(t)\);
		\item an extensive variable \(\Phi = \Phi(t)\) with its intensive counterpart
		      \(\phi = \phi(t, \vec{x})\)  defined inside the material volume.
	\end{enumerate}
	The total derivative of \(\Phi\) with respect to time can be formulated as:
	\begin{equation}
    \odv{\Phi}{t} = \odv{}{t} \int\limits_{\Omega(t)} \phi \d \Omega
		= \int\limits_{\Omega(t)} \mdv{\phi}{t} + \phi \div{\vec{v}} \d \Omega
	\end{equation}
\end{theorem}

The theorem in proved below. We can rephrase it slightly and use it 
\[
  \odv{}{t}\int\limits_\Omega \rho E \d \Omega 
  = \int\limits_\Omega \rho \mdv{E}{t} \d \Omega.
\]

All the expressions are integrals in \(\Omega\). We can drop the integration as the only way they 
are equal is if the integrated \say{things} are equal to each other. So we end up having
\[
  \rho\mdv{E}{t} = -\pdv{pu_i}{x_i} + \pdv{u_i\sigma_{ij}}{x_j} - \pdv{q_i}{x_i} + \rho\ab(u_i f_i + Q),
\]
which is the equation for total energy.

\section{Internal energy equation}
We subtract the kinetic energy equation from the total energy equation. Make
the exercise yourself. Expand the terms \(\pdv{p u_i}{x_i}\) and
\(\pdv{\sigma_{ij}u_i}{x_j}\). The final equation is
\[
  \rho\mdv{e}{t} = -p\pdv{u_i}{x_i} + \sigma_{ij}\pdv{u_i}{x_j} - \pdv{q_i}{x_i} + \rho Q.
\]
Note no work of body forces.

\subsection{Enthalpy equation}
Enthalpy is 
\[
  h = e + \frac{p}{\rho} \rightarrow e = h - \frac{p}{\rho}
\]
so 
\[
  \rho\mdv{e}{t} = \pdv{\rho (h - p/\rho)}{t} + \pdv{\rho u_i (h - p/\rho)}{x_i} =
  \rho\mdv{h}{t} - \mdv{p}{t}
\]
If we substitute into internal energy we get
\[
  \rho\mdv{h}{t} = \mdv{p}{t} -p\pdv{u_i}{x_i} + \sigma_{ij}\pdv{u_i}{x_j} - \pdv{q_i}{x_i} + \rho Q.
\]
\subsection{Total enthalpy}
We have to add the kinetic energy equation. Interesting stuff is happening only on right side.
Summing both equations
\[
  \begin{split}
    \rho\mdv{H}{t} &= \underbrace{-\pdv{p}{t} + \pdv{p u_i}{x_i}}_{\mdv{p}/{t}} 
    -p\pdv{ u_i}{x_i} - \pdv{p}{x_i}u_i + u_i\pdv{\sigma_{ij}}{x_j} + \sigma_{ij}\pdv{u_i}{x_j} - \pdv{q_i}{x_i}
    + \rho\ab(u_i f_i + Q)\\
    &= \pdv{p}{t} + \pdv{u_i \sigma_{ij}}{x_j} - \pdv{q_i}{x_i}
    + \rho\ab(u_i f_i + Q)
  \end{split}
\]


\section{Proof or Reynolds transform theorem (extra material)}


\paragraph{Material volume}
Special integration region \(\Omega\) defined as a parcel of fluid
with such boundaries that there is no flux of selected intensive property 
through them.

\paragraph{Index notation and Levi-Civita symbol}
Index notation, or Einstein notation simplifies greatly multi variable 
calculus and relies on two rules:
\begin{enumerate}
  \item Any vector variable can be represented as symbol with an index.
    For example, velocity \(\vec{v}\), a vector having generally three 
    components, can be simply written as \(v_i\). Index \(i\) is the so 
    called \textbf{free index}. The letter does not matter, i.e.\ \(i\) can
    be replaced with \(\alpha\) or whatever.
  \item Repeated index implies summation over full range of this index. 
    This means that, if a term like \(a_i b_j c_j\) is encountered, it can
    be expanded into
    \[
      a_i b_j c_j = a_i b_1 c_1 + a_i b_2 c_2 + a_i b_3 c_3 =
      \sum\limits_{j=1,2,3} a_i b_j c_j.
    \]
    Therefore, we can say that the sum sign \(\Sigma\) is "implicit". The index
    \(i\) remained a free index in the above expression. This means that a
    version for each spatial dimension can be obtained by setting \(i\) to 
    a chosen index.
\end{enumerate}
The Levi-Civita symbol, \(\levicivita{i}{j}{k}\), is a helpful permutation
symbol defined in a following way:
\begin{equation*}
\varepsilon_{ijk} = \begin{cases}
         +1 & \text{if } (i,j,k) \text{ is } (1,2,3), (2,3,1), \text{ or } (3,1,2), \\
         -1 & \text{if } (i,j,k) \text{ is } (3,2,1), (1,3,2), \text{ or } (2,1,3), \\
    \;\;\,0 & \text{if } i = j, \text{ or } j = k, \text{ or } k = i
\end{cases}
\end{equation*}

\paragraph{Transforming material volume to reference configuration}
The integration region is usually the function of time, \(\Omega = \Omega(t)\).
The reference configuration of \(\Omega(t)\), denoted as \(\Omega^0\), is not 
time independent. Assuming that transformation \(\Omega^0 \rightarrow \Omega\)
is invertable, we can change the integration region to the reference
configuration by writing:
\begin{equation}
  d\Omega(t) = J(t)d\Omega^0,
  \label{eq::jacobian_relation}
\end{equation}
where \(J\) denotes a Jacobian of the transformation.

\paragraph{Defining the Jacobian}
The trajectory of a point \(\vec{x}\) laying inside 
the parcel \(\Omega\) is a function of time and the initial 
configuration \(\vec{x}^0\):
\begin{equation}
  \vec{x} = \vec{x}(t, \vec{x}^0)
  \label{eq::trajectory}
\end{equation}
In the above, we have assumed that for a given point it is possible 
to find its reference configuration, that is inverse function 
\(\vec{x}^0 = \vec{x}^0(t, \vec{x})\) exists.
The necessary and sufficient condition for invertability is the 
non-vanishing Jacobian, which can be given as:
\begin{equation}\label{eq::jacobian}
  J = \det\pdv{x_i}{x^0_j} 
  = \levicivita{i}{j}{k} \pdv{x_1}{x^0_i}\pdv{x_2}{x^0_j}\pdv{x_3}{x^0_k}
\end{equation}

\paragraph{Time derivative of intensive variable, substantial derivative}
Total time derivative of an intensive property \(d \phi / d t\) 
(also called substantial derivative, denoted \(D\phi/Dt\)) can be formulated 
in the following way. We assume that trajectory of each point inside a fluid can
be described as in \autoref{eq::trajectory} and 
we use the chain rule:
\begin{equation*}
  \mdv{\phi}{t} \equiv \odv{}{\phi}{t}  
  =  \odv{}{t}f\ab(\vec{x}\ab(t,\vec{x}^0))
  = \pdv{\phi}{t} + \odv{x_i}{t} \pdv{\phi}{x_i}
  = \pdv{\phi}{t} + \ab(\odv{\vec{x}}{t} \cdot \grad )\phi
\end{equation*}
Since the time derivative of the trajectory \(d\vec{x} / dt\) is equal 
to the fluids velocity \(\vec{v}\) by definition, the above reduces to:
\begin{equation}
  \mdv{\phi}{t} \equiv \odv{\phi}{t} 
  = \pdv{\phi}{t} + (\vec{v} \cdot \grad) \phi
\end{equation}


\paragraph{Material derivative of the Jacobian}
The Jacobian arises most often while changing the integration region 
from \(\Omega(t)\) to \(\Omega^0\) according to \autoref{eq::jacobian_relation}.
The closed expression for the time derivative of the Jacobian must be found 
for the proof or RTT.
Starting from the definition of the Jacobian determinant~\autoref{eq::jacobian},
and distributing the derivative:
\begin{align*}
  \odv{}{J}{t} = \odv{}{t} \ab(\levicivita{i}{j}{k} 
  \pdv{x_1}{x^0_i} \pdv{x_2}{x^0_j} \pdv{x_3}{x^0_k})
  &= \levicivita{i}{j}{k} \odv{}{t} \ab( \pdv{x_1}{x^0_i} ) 
    \pdv{x_2}{x^0_j} \pdv{x_3}{x^0_k} \\
  &+ \levicivita{i}{j}{k} \pdv{x_1}{x^0_i} 
    \odv{}{t} \ab( \pdv{x_2}{x^0_j})  \pdv{x_3}{x^0_k} 
  + \levicivita{i}{j}{k} \pdv{x_1}{x^0_i} 
    \pdv{x_2}{x^0_j} \odv{}{t} \ab( \pdv{x_3}{x^0_k} )  
\end{align*}
In the above, differentiation with respect to time and reference
configuration coordinates commute (the order of differentiation 
can be freely changed), resulting in velocity components:
\begin{equation*}
  \odv{}{t} \ab( \pdv{x_\alpha}{x^0_\beta} )
  = \pdv{v_\alpha}{x^0_\beta}
  = \pdv{v_\alpha}{x_\gamma} \pdv{x_\gamma}{x^0_\beta}
\end{equation*}
We also insert additional differentiation for the sake of 
further derivation.

By inserting the above into the formula for the Jacobian derivative 
and simplifying, we get final formula for the \(\d  J/\d  t\).
Taking only one of the terms as an example and expanding summation 
over \(\gamma\):
\begin{equation*}
  \begin{split}
    & \levicivita{i}{j}{k} \odv{}{t} \ab( \pdv{x_1}{x^0_i} ) 
      \pdv{x_2}{x^0_j} \pdv{x_3}{x^0_k}
    = \pdv{v_1}{x_\gamma} \levicivita{i}{j}{k} \pdv{x_\gamma}{x^0_i}\pdv{x_2}{x^0_j}
      \pdv{x_3}{x^0_k} =\\
    & \pdv{v_1}{x_1} \levicivita{i}{j}{k} \pdv{x_1}{x^0_i}\pdv{x_2}{x^0_j} \pdv{x_3}{x^0_k} + 
    \pdv{v_1}{x_2} \levicivita{i}{j}{k} \pdv{x_2}{x^0_i}\pdv{x_2}{x^0_j} \pdv{x_3}{x^0_k} +
    \pdv{v_1}{x_3} \levicivita{i}{j}{k} \pdv{x_3}{x^0_i}\pdv{x_2}{x^0_j} \pdv{x_3}{x^0_k}
  \end{split}
\end{equation*}
Both terms on the right with repeated index in numerators will be equal to 
zero because of the definition of \(\levicivita{i}{j}{k}\). They will be 
repeated two times, once positive and once negative, for example:
\[
  \pdv{v_1}{x_2} \levicivita{i}{j}{k} 
  \pdv{x_2}{x^0_i}\pdv{x_2}{x^0_j} \pdv{x_3}{x^0_k} = \dots
  \pdv{v_1}{x_2} 
  \pdv{x_2}{x^0_1}\pdv{x_2}{x^0_2} \pdv{x_3}{x^0_3} + \dots 
  - \pdv{v_1}{x_2} 
  \pdv{x_2}{x^0_2}\pdv{x_2}{x^0_1} \pdv{x_3}{x^0_3} + \dots
\]
The terms on the right cancel each other out. In the above, we shown only 
\((i,j,k) = (1,2,3)\) and \((2,1,3)\).
Thanks to that, and using the original Jacobian definition, 
i.e.~\autoref{eq::jacobian}, we can simplify first of the terms: 
\begin{equation*}
  \levicivita{i}{j}{k} \odv{}{t} \ab( \pdv{x_1}{x^0_i} ) \pdv{x_2}{x^0_j} 
  \pdv{x_3}{x^0_k}
  = \pdv{v_1}{x_1} \levicivita{i}{j}{k} \pdv{x_1}{x^0_i}\pdv{x_2}{x^0_j} 
  \pdv{x_3}{x^0_k}
  = \pdv{v_1}{x_1} J
\end{equation*}
Treating indices \(2\) and \(3\) the same way, the final formula for the 
time derivative of the Jacobian looks as follows:
\begin{equation}
  \odv{J}{t}  = \pdv{v_i}{x_i} J = \ab(\div\vec{v})J
  \label{eq::ddt_J}
\end{equation}

\begin{proof}
	The integral of \(\phi\) can be expressed in reference configuration 
  \(\Omega^0\), assuming that \(\Omega(t) = J(t)\Omega^0\):
	\begin{equation}
    \odv{}{t} \int\limits_{\Omega(t)} \phi(t, \vec{x}) \d{\Omega} = 
    \odv{}{t} \int\limits_{\Omega^0} \phi(t, \vec{x^0})J(t) \d{\Omega^0} 
	\end{equation}
	Since \(\Omega_0\) is time independent, the derivative can be moved
  under the integral:
	\begin{equation}
    \odv{}{t} \int\limits_{\Omega(t)} \phi(t, \vec{x}) \d{\Omega} = 
    \int\limits_{\Omega^0} \odv{}{t}[\phi(t, \vec{x^0})J(t)] \d{\Omega^0} =
    \int\limits_{\Omega^0} \mdv{\phi}{t}J + \phi\odv{}{J}{t} \d{\Omega^0} 
	\end{equation}
  Note the presence of substantial derivative; the intensive quantity is
  transported by the fluid.
	Factoring the Jacobian out of the terms under the integral, we can return
  to the original integration region
	\begin{equation}
    \int\limits_{\Omega^0} \mdv{\phi}{t}J + \phi\odv{}{J}{t} \d{\Omega^0} = 
    \int\limits_{\Omega^0} \ab[\mdv{\phi}{t} + J^{-1}\phi\odv{J}{t}]J \d{\Omega^0} = 
    \int\limits_{\Omega(t)} \mdv{\phi}{t} + J^{-1}\phi\odv{J}{t} \d{\Omega}.
	\end{equation}
  The final formula is obtained by plugging in the time derivative of the 
  Jacobian~\autoref{eq::ddt_J}:
  \begin{equation*}
    \int\limits_{\Omega(t)} \mdv{\phi}{t} + J^{-1}\phi\odv{}{J}{t} \d{\Omega}
    = \int\limits_{\Omega(t)} \mdv{\phi}{t} + \phi \div\vec{v} \d{\Omega}
  \end{equation*}
\end{proof}




\end{document}
